\documentclass[11pt]{article}

% ===================== PACKAGES =====================
\usepackage[a4paper,margin=1in]{geometry}
\usepackage{amsmath,amssymb}
\usepackage{enumitem}
\usepackage{fancyhdr}
\usepackage{xcolor} 

% ===================== HEADER & FOOTER =====================
\pagestyle{fancy}
\fancyhf{}
\lhead{CSE 400: Fundamentals of Probability in Computing}
\rhead{Lecture-11 Scribe Submission}
\cfoot{\thepage}

% ===================== CUSTOM COMMANDS =====================
\newcommand{\solution}{
    \vspace{0.5cm}
    \noindent\textbf{\textcolor{blue}{Solution:}}
    \vspace{0.2cm}
    
}

% ===================== TITLE =====================
\title{
    \normalsize School of Engineering and Applied Science (SEAS), Ahmedabad University \\
    \vspace{0.2cm}
    \textbf{CSE 400: Fundamentals of Probability in Computing}\\
    \Large Lecture-11 Scribe Submission
}
\author{} 
\date{}

\begin{document}
\maketitle

\vspace{-2cm}
\begin{center}
    \begin{tabular}{ll}
        \textbf{Name:} & {Jashvi Shah \hspace{2.5in}} \\ [1.5ex]
        \textbf{Enrollment No:} & {AU2440026 \hspace{2.5in}} \\ [1.5ex]
        \textbf{Email:} & {jashvi.s@ahduni.edu.in \hspace{2.1in}} \\ [1.5ex]
    \end{tabular}
\end{center}

\hrule
\vspace{0.5cm}

\section*{CSE400 – Lecture 11 Scribe}

\date{Instructor: Dhaval Patel, PhD \\ Date: February 10, 2026}

\begin{document}

\maketitle

\section{Lecture Structure and Objective}

\subsection{Outline of the Lecture}

The lecture consisted of three main components:

\begin{enumerate}
    \item Transformation of Random Variables
    \begin{itemize}
        \item Learning transformation techniques for random variables.
    \end{itemize}
    
    \item Function of Two Random Variables
    \begin{itemize}
        \item Understanding joint transformations and derived distributions.
    \end{itemize}
    
    \item Illustrative Example
    \begin{itemize}
        \item Detailed derivation for the case:
        \[
        Z = X + Y
        \]
    \end{itemize}
\end{enumerate}

\section{Transformation of Random Variables}

\subsection{Definition and Notation}

Let $X$ denote an original random variable.

Let $Z$ denote a new random variable obtained through transformation of $X$.

Thus, a transformation is represented as:
\[
Z = g(X)
\]

\subsection{Objective of Transformation}

Transformation techniques are used to:

\begin{itemize}
    \item Derive new random variables from existing ones.
    \item Study how the properties of the original variable affect the transformed variable.
\end{itemize}

\subsection{Reasoning and Logical Development}

Step-wise reasoning presented in lecture:

\begin{enumerate}
    \item Begin with a known random variable.
    \item Apply a mathematical function to the variable.
    \item The output of the function becomes a new random variable.
    \item Transformation techniques help determine the behavior and characteristics of the new variable.
\end{enumerate}

\subsection{Assumptions and Conditions}

\begin{itemize}
    \item The transformation is defined using a function applied to an existing random variable.
    \item The resulting value is treated as a random variable.
\end{itemize}

\section{Function of Two Random Variables}

\subsection{Definition and Notation}

Let $X$ and $Y$ denote two random variables.

Let $Z$ denote a new random variable obtained from both variables.

The transformation is represented as:
\[
Z = g(X,Y)
\]

\subsection{Concept of Joint Transformation}

Joint transformation involves:

\begin{itemize}
    \item Combining two random variables.
    \item Producing a new random variable using a function of both variables.
    \item Studying derived distributions based on the joint behavior of the original variables.
\end{itemize}

\subsection{Logical Step-by-Step Reasoning}

\begin{enumerate}
    \item Start with two random variables.
    \item Define a function involving both variables.
    \item Apply the function to generate a new random variable.
    \item Analyze the derived distribution of the new variable.
\end{enumerate}

\subsection{Result Emphasized in Lecture}

\begin{itemize}
    \item Functions of multiple random variables lead to derived distributions.
    \item Joint transformation extends the idea of transformation from one variable to multiple variables.
\end{itemize}

\section{Illustrative Example: $Z = X + Y$}

\subsection{Statement of the Example}

The lecture provided a detailed derivation example using:
\[
Z = X + Y
\]

\subsection{Objective of the Example}

\begin{itemize}
    \item To demonstrate practical application of joint transformation.
    \item To illustrate how a new random variable is formed from two variables.
\end{itemize}

\subsection{Step-by-Step Reasoning Provided}

\begin{enumerate}
    \item Consider two random variables $X$ and $Y$.
    \item Define a new variable $Z$ as the sum of the two variables.
    \item The value of $Z$ depends on the combined values of $X$ and $Y$.
    \item The derived distribution of $Z$ is determined from the joint behavior of $X$ and $Y$.
\end{enumerate}

\subsection{Key Insight from Example}

\begin{itemize}
    \item Summation of two random variables is a special case of transformation involving two variables.
    \item The example demonstrates how transformation techniques are applied to derive new distributions.
\end{itemize}

\section{Logical Flow of Ideas Across the Lecture}

The lecture followed the progression:

\begin{enumerate}
    \item Introduction to transformation of a single random variable.
    \item Extension of transformation techniques to functions involving two random variables.
    \item Application of concepts using the summation example.
\end{enumerate}

\section{Key Takeaways}

\begin{itemize}
    \item Transformation methods allow creation and analysis of new random variables.
    \item Joint transformations consider relationships between multiple variables.
    \item Derived distributions describe outcomes resulting from transformations.
    \item The example $Z = X + Y$ demonstrates practical implementation of transformation concepts.
\end{itemize}

\end{document}
